% GENERATED FILE. Edit canonical lessons, then run npm run generate:books.
% canonical-source-hash: fnv1a64:a5417dfe5dce8036
% canonical-lessons: MW-C177-pharno, MW-C177-batan, MW-C177-takno, MW-C177-theglo, MW-C177-pallo

\chapter{Torn and Mended}
\label{ch:mw-torn-and-mended}
\hlchaptermodality{177}

\begin{chapteropening}
\textbf{By the end of this chapter:} \emph{I can say that something tore and how it was mended.}

Complete the last lesson of chapter 177: I can say that something tore and how it was mended.
\end{chapteropening}

\section[\texorpdfstring{\mw{फाड़णो} (phāṛṇo)}{फाड़णो (phāṛṇo)}]{\mw{फाड़णो} (phāṛṇo) — to tear}

\label{lesson:MW-C177-pharno}



\begin{quote}
Before the new one: say the Marwadi for a door chain, then the Marwadi for to hang up.
\end{quote}

\subsection*{You'll want to know: \mw{फाड़णो}}
\textbf{\mw{फाड़णो}} — \emph{phāṛṇo} — \textquotedblleft{}to tear\textquotedblright{}.

\begin{quote}
\mw{छोरो} \mw{चिट्ठी} \mw{फाड़ै} \mw{है।}
\end{quote}

\begin{quote}
\emph{chhoro chiṭṭhī phāṛai hai.}
\end{quote}

\begin{quote}
The boy tears up the letter.
\end{quote}

Past (I …): \mw{म्हैं} \mw{चिट्ठी} \mw{फाड़ी}, \emph{mhaĩ chiṭṭhī phāṛī}.

\begin{cousinweb}
\textbf{\mw{फाड़णो}} (\emph{phāṛṇo}) is inherited from Sanskrit \emph{sphāṭayati}, \textquotedblleft{}causes to split\textquotedblright{}; CDIAL lists Old Marwari \emph{phāḍaï}, \textquotedblleft{}to split\textquotedblright{} (CDIAL 13825).
\end{cousinweb}

\begin{grammarlens}[title={What you've built}]
One more word for this chapter.
\end{grammarlens}

\subsection*{Your turn}
\begin{itemize}
\raggedright
  \item \emph{Say it:} \emph{phāṛṇo}
  \item \emph{Say it:} \emph{phāṛṇo}, once more
  \item \emph{Say it:} \emph{chhoro chiṭṭhī phāṛai hai.} — The boy tears up the letter.
  \item \emph{Recall:} say the Marwadi for a door chain, then the Marwadi for to hang up, then say \emph{phāṛṇo} again
\end{itemize}

\begin{tcolorbox}[breakable,colback=teal!4,colframe=teal!35!black,title={Before you move on}]
What does \mw{फाड़णो} mean? (\textquotedblleft{}To tear\textquotedblright{}.) Say it once more.
\end{tcolorbox}
\section[\texorpdfstring{\mw{बटण} (baṭaṇ)}{बटण (baṭaṇ)}]{\mw{बटण} (baṭaṇ) — a button}

\label{lesson:MW-C177-batan}



\begin{quote}
Before the new one: say the Marwadi for to hang up, then the Marwadi for to tear.
\end{quote}

\subsection*{You'll want to know: \mw{बटण}}
\textbf{\mw{बटण}} — \emph{baṭaṇ} — \textquotedblleft{}a button\textquotedblright{}.

\begin{quote}
\mw{म्हारी} \mw{कमीज} \mw{माथै} \mw{नुवो} \mw{बटण} \mw{है।}
\end{quote}

\begin{quote}
\emph{mhārī kamīj māthai nuvo baṭaṇ hai.}
\end{quote}

\begin{quote}
There's a new button on my shirt.
\end{quote}

A masculine noun.

\begin{cousinweb}
\textbf{\mw{बटण}} (\emph{baṭaṇ}) is borrowed from English \emph{button}.
\end{cousinweb}

\begin{grammarlens}[title={What you've built}]
One more word for this chapter.
\end{grammarlens}

\subsection*{Your turn}
\begin{itemize}
\raggedright
  \item \emph{Say it:} \emph{baṭaṇ}
  \item \emph{Say it:} \emph{baṭaṇ}, once more
  \item \emph{Say it:} \emph{mhārī kamīj māthai nuvo baṭaṇ hai.} — There's a new button on my shirt.
  \item \emph{Recall:} say the Marwadi for to hang up, then the Marwadi for to tear, then say \emph{baṭaṇ} again
\end{itemize}

\begin{tcolorbox}[breakable,colback=teal!4,colframe=teal!35!black,title={Before you move on}]
What does \mw{बटण} mean? (\textquotedblleft{}A button\textquotedblright{}.) Say it once more.
\end{tcolorbox}
\section[\texorpdfstring{\mw{टांकणो} (ṭā̃kṇo)}{टांकणो (ṭā̃kṇo)}]{\mw{टांकणो} (ṭā̃kṇo) — to stitch on}

\label{lesson:MW-C177-takno}



\begin{quote}
Before the new one: say the Marwadi for to tear, then the Marwadi for a button.
\end{quote}

\subsection*{You'll want to know: \mw{टांकणो}}
\textbf{\mw{टांकणो}} — \emph{ṭā̃kṇo} — \textquotedblleft{}to stitch on\textquotedblright{}.

\begin{quote}
\mw{मां} \mw{कमीज} \mw{माथै} \mw{बटण} \mw{टांकै} \mw{है।}
\end{quote}

\begin{quote}
\emph{mā̃ kamīj māthai baṭaṇ ṭā̃kai hai.}
\end{quote}

\begin{quote}
Mother stitches a button onto the shirt.
\end{quote}

Past (I …): \mw{म्हैं} \mw{बटण} \mw{टांक्यो}, \emph{mhaĩ baṭaṇ ṭā̃kyo}.

\begin{cousinweb}
\textbf{\mw{टांकणो}} (\emph{ṭā̃kṇo}) is inherited from Sanskrit \emph{ṭaṅkati}, \textquotedblleft{}ties\textquotedblright{}; CDIAL lists Hindi \emph{ṭā̃knā}, \textquotedblleft{}to stitch, join\textquotedblright{} (CDIAL 5432). CDIAL gives no Marwari form in this entry.
\end{cousinweb}

\begin{grammarlens}[title={What you've built}]
One more word for this chapter.
\end{grammarlens}

\subsection*{Your turn}
\begin{itemize}
\raggedright
  \item \emph{Say it:} \emph{ṭā̃kṇo}
  \item \emph{Say it:} \emph{ṭā̃kṇo}, once more
  \item \emph{Say it:} \emph{mā̃ kamīj māthai baṭaṇ ṭā̃kai hai.} — Mother stitches a button onto the shirt.
  \item \emph{Recall:} say the Marwadi for to tear, then the Marwadi for a button, then say \emph{ṭā̃kṇo} again
\end{itemize}

\begin{tcolorbox}[breakable,colback=teal!4,colframe=teal!35!black,title={Before you move on}]
What does \mw{टांकणो} mean? (\textquotedblleft{}To stitch on\textquotedblright{}.) Say it once more.
\end{tcolorbox}
\section[\texorpdfstring{\mw{थेगळो} (thegḷo)}{थेगळो (thegḷo)}]{\mw{थेगळो} (thegḷo) — a patch (of cloth)}

\label{lesson:MW-C177-theglo}



\begin{quote}
Before the new one: say the Marwadi for a button, then the Marwadi for to stitch on.
\end{quote}

\subsection*{You'll want to know: \mw{थेगळो}}
\textbf{\mw{थेगळो}} — \emph{thegḷo} — \textquotedblleft{}a patch (of cloth)\textquotedblright{}.

\begin{quote}
\mw{मां} \mw{थेगळो} \mw{सीवै} \mw{है।}
\end{quote}

\begin{quote}
\emph{mā̃ thegḷo sīvai hai.}
\end{quote}

\begin{quote}
Mother sews on a patch.
\end{quote}

A masculine noun.

\begin{cousinweb}
\textbf{\mw{थेगळो}} (\emph{thegḷo}) comes from a reconstructed \emph{*thigga}, \textquotedblleft{}a rag\textquotedblright{}; CDIAL lists Hindi \emph{theglī}, \textquotedblleft{}a patch of cloth\textquotedblright{} (CDIAL 6096). CDIAL gives no Marwari form in this entry.
\end{cousinweb}

\begin{grammarlens}[title={What you've built}]
One more word for this chapter.
\end{grammarlens}

\subsection*{Your turn}
\begin{itemize}
\raggedright
  \item \emph{Say it:} \emph{thegḷo}
  \item \emph{Say it:} \emph{thegḷo}, once more
  \item \emph{Say it:} \emph{mā̃ thegḷo sīvai hai.} — Mother sews on a patch.
  \item \emph{Recall:} say the Marwadi for a button, then the Marwadi for to stitch on, then say \emph{thegḷo} again
\end{itemize}

\begin{tcolorbox}[breakable,colback=teal!4,colframe=teal!35!black,title={Before you move on}]
What does \mw{थेगळो} mean? (\textquotedblleft{}A patch (of cloth)\textquotedblright{}.) Say it once more.
\end{tcolorbox}
\section[\texorpdfstring{\mw{पल्लो} (pallo)}{पल्लो (pallo)}]{\mw{पल्लो} (pallo) — the end of a sari; the loose end of a head-cloth}

\label{lesson:MW-C177-pallo}



\begin{quote}
Before the new one: say the Marwadi for to stitch on, then the Marwadi for a patch.
\end{quote}

\subsection*{You'll want to know: \mw{पल्लो}}
\textbf{\mw{पल्लो}} — \emph{pallo} — \textquotedblleft{}the end of a sari; the loose end of a head-cloth\textquotedblright{}.

\begin{quote}
\mw{म्हारो} \mw{पल्लो} \mw{गीलो} \mw{है।}
\end{quote}

\begin{quote}
\emph{mhāro pallo gīlo hai.}
\end{quote}

\begin{quote}
The end of my head-cloth is wet.
\end{quote}

A masculine noun.

\begin{cousinweb}
\textbf{\mw{पल्लो}} (\emph{pallo}) is inherited from Sanskrit \emph{pallava}, \textquotedblleft{}a strip of cloth\textquotedblright{}; CDIAL lists Old Marwari \emph{palā}, \textquotedblleft{}the corners of a garment\textquotedblright{} (CDIAL 7970). It is a different word from \emph{pallava}, \textquotedblleft{}a sprout\textquotedblright{}.
\end{cousinweb}

\begin{grammarlens}[title={What you've built}]
One more word for this chapter.
\end{grammarlens}

\subsection*{Your turn}
\begin{itemize}
\raggedright
  \item \emph{Say it:} \emph{pallo}
  \item \emph{Say it:} \emph{pallo}, once more
  \item \emph{Say it:} \emph{mhāro pallo gīlo hai.} — The end of my head-cloth is wet.
  \item \emph{Recall:} say the Marwadi for to stitch on, then the Marwadi for a patch, then say \emph{pallo} again
\end{itemize}

\begin{tcolorbox}[breakable,colback=teal!4,colframe=teal!35!black,title={Before you move on}]
What does \mw{पल्लो} mean? (\textquotedblleft{}The end of a sari; the loose end of a head-cloth\textquotedblright{}.) Say it once more.
\end{tcolorbox}
